% Abstract — problem → gap → method → results (research-story order)
\begin{abstract}
Deploying large language model (LLM) agents requires selecting among heterogeneous agent
strategies, including direct prompting, chain-of-thought reasoning, tool-augmented ReAct, and
multi-agent coordination.
Because no single strategy is optimal across queries, effective deployment depends on routing
queries to appropriate strategies before execution.
Existing routing approaches primarily rely on semantic query representations and hard
winner-take-all supervision, which may overlook procedural requirements and discard information
when multiple strategies succeed at different costs.

We propose \textbf{Adaptive Agent Routing (AAR)}, a pre-inference routing framework that combines
planner-derived procedural complexity with semantic query representations.
AAR first estimates query complexity from an LLM-generated procedure directed acyclic graph
(DAG), summarized into a five-dimensional complexity vector.
This signal is fused with semantic embeddings to form a complexity--semantic routing
representation.
To learn routing policies, AAR employs \textbf{cost-soft supervision}, which distributes
training mass across successful strategies according to execution cost rather than enforcing a
single hard target.

Experiments on a held-out routing benchmark ($n{=}100$ test; $\lambda{=}25$) show that
\textbf{cost-soft supervision} substantially reduces utility regret (agent regret
$0.361\to 0.192$; executed exact match 63\%$\to$82\%).
The lowest total regret ($0.198$) is achieved when cost-soft training is combined with the
complexity--semantic representation, outperforming embedding-only, complexity-only, and
heuristic-plus-embedding baselines.
Representation ablations indicate that procedural complexity alone is insufficient for routing
but provides complementary information when combined with semantic features.
Under transfer on executed benchmark workloads ($N{=}965$), AAR attains competitive performance
(48.1\% pooled exact match vs.\ 50.3\% for a heuristic baseline); gains are less consistent
than on the routing benchmark, where oracle multi-strategy supervision aligns with the routing
objective.
\end{abstract}
